%%% FIELD profile-negative-statement
Fix \(1<s\le2\), \(M\ge2\), and \(v\in(0,1/64)\).  The asserted
uniform localization inequality for
\((\widehat c_{\mathrm{prof}},\mathsf F_{\mathrm{prof}})\) is false on
\(\mathcal P_n(s,M,v,\kappa_n)\).  More precisely, for every
\[
 0<\varepsilon<\frac{1}{2(2+1/s)},
 \qquad \kappa_n=\frac18n^{-\varepsilon},
 \qquad H_n=\left(\frac{\kappa_n}{2^{12}M}\right)^{1/s},
 \qquad m=\lfloor n/4\rfloor,
\]
there is, for all sufficiently large \(n\), a law
\(P_n\in\mathcal P_n(s,M,v,\kappa_n)\) with \(c_{P_n}=1/2\) such that
\[
 \frac{T_n}{\log(e/\kappa_n)}\longrightarrow\infty,
 \qquad \frac{16}{n}\le\frac{H_n}{8},
\]
and
\[
 P_n\{\mathsf F_{\mathrm{prof}}=1\}
 +P_n\!\left(\left\{|\widehat c_{\mathrm{prof}}-c_{P_n}|>
                    H_n/8\right\}
       \cap\{\mathsf F_{\mathrm{prof}}=0\}\right)
 \longrightarrow1.
\]
Consequently, there are no positive constants
\(A_{\mathrm{prof}},a_{\mathrm{prof}},C_{\mathrm{prof}}\), depending
only on \(s,M\), such that, simultaneously for every
\(P\in\mathcal P_n(s,M,v,\kappa_n)\), every index satisfying
\(T_n\ge C_{\mathrm{prof}}\log(e/\kappa_n)\), and every
\(16/n\le u\le H_n\),
\[
 P\{\mathsf F_{\mathrm{prof}}=1\}
 +P\!\left(\{|\widehat c_{\mathrm{prof}}-c_P|>u\}
       \cap\{\mathsf F_{\mathrm{prof}}=0\}\right)
 \le A_{\mathrm{prof}}\exp\{-a_{\mathrm{prof}}m\kappa_n^2u\}.
\]

%%% FIELD profile-negative-proof
Fix \(s\in(1,2]\), \(M\ge2\), \(v\in(0,1/64)\), and choose
\[
 0<\varepsilon<\frac{1}{2(2+1/s)},
 \qquad
 \kappa_n=\frac18 n^{-\varepsilon},
 \qquad
 H_n=\left(\frac{\kappa_n}{2^{12}M}\right)^{1/s},
 \qquad c=\frac12,
 \qquad m=\lfloor n/4\rfloor.
 \tag{1}
\]
Write \(q_n=\kappa_n/H_n\).  Since \(s>1\),
\[
 q_n=(2^{12}M)^{1/s}\kappa_n^{\,1-1/s}\longrightarrow0.
 \tag{2}
\]

We first construct the full-data law.  Let \(X\) be uniform on
\([0,1]\).  Conditional on \(X=x\), assign the monotone response type
\(R\) the probabilities
\[
 P_n(R=A\mid X=x)=\frac14,
 \qquad
 P_n(R=C\mid X=x)=w_n(x),
 \qquad
 P_n(R=N\mid X=x)=\frac34-w_n(x),
 \tag{3}
\]
where
\[
 w_n(x)=
 \begin{cases}
  \kappa_n,&x<c,\\
  \kappa_n+q_n(x-c),&x\ge c.
 \end{cases}
 \tag{4}
\]
By (2), \(0\le w_n(x)\le\kappa_n+q_n/2<1/4\) for all sufficiently
large \(n\); hence the three masses in (3) are nonnegative and sum to
one.  Given \((X,R)\), let \(U\) be Bernoulli with success probability
\(1/2\), independently of \((X,R)\), and set
\[
 Y(0)=Y(1)=U.
 \tag{5}
\]
Use the frozen threshold assignment \(Z=\mathbf 1\{X\ge c\}\), the
frozen monotone treatment rule \(D=D(Z)\), and independent copies
across units.

We verify every condition defining the model class.  From (3)--(4),
\[
 E_{P_n}[D(0)\mid X=x]=\frac14,
 \qquad
 E_{P_n}[D(1)\mid X=x]=\frac14+w_n(x).
 \tag{6}
\]
Both functions in (6) are continuous at \(c\), because the two
one-sided values of \(w_n\) there are \(\kappa_n\).  By (5),
\(E_{P_n}[Y(D(z))\mid X=x]=1/2\) for both \(z\in\{0,1\}\).  Thus
\(\mathrm{ass:potential\mbox{-}reduced\mbox{-}form\mbox{-}continuity}\)
holds.  The observable treatment regressions are
\[
 \mu_{D,P_n}^{-}(x)=\frac14,
 \qquad
 \mu_{D,P_n}^{+}(x)=\frac14+\kappa_n+q_n(x-c).
 \tag{7}
\]
The left side has a constant extension, and the displayed affine
formula extends the right side to \([0,1]\).  For all sufficiently
large \(n\), these extensions take values in \([0,1]\), have uniformly
bounded suprema, have derivative suprema at most \(q_n\), and have
zero derivative H\"older seminorm.  Equation (2) and \(M\ge2\)
therefore place them in the radius-\(M\), order-\(s\) H\"older ball
required by
\(\mathrm{ass:observable\mbox{-}holder\mbox{-}extensions}\).  The two
observable outcome regressions are the constant \(1/2\), so they obey
the same condition.  Moreover,
\[
 d_{P_n}=\mu_{D,P_n}^{+}(c)-\mu_{D,P_n}^{-}(c)=\kappa_n,
 \tag{8}
\]
which verifies \(\mathrm{ass:local\mbox{-}compliance\mbox{-}strength}\).
Finally, (5) and the independence of \(U\) give, on both sides
\(\ell\),
\[
 \operatorname{Var}_{P_n}(Y\mid X=x,\ell)=\frac14>v,
 \qquad
 \operatorname{Cov}_{P_n}(Y,D\mid X=x,\ell)=0.
 \tag{9}
\]
Thus the variance floor and correlation-separation conditions hold.
The sampling is independent, the design is uniform, and \(c=1/2\) is
interior.  Equations (3), (5)--(9), together with those three facts,
prove
\[
 P_n\in\mathcal P_n(s,M,v,\kappa_n)
 \quad\text{for all sufficiently large }n.
 \tag{10}
\]

We next evaluate exactly the population criterion induced by the
moving-window profile.  For a candidate \(t\) satisfying
\(|t-c|\le H_n\), put
\[
 z=\frac{x-t}{H_n},
 \qquad
 r=\frac{c-t}{H_n},
 \qquad
 f_r(z)=\mathbf 1\{z\ge r\}(1+z-r).
 \tag{11}
\]
After subtracting the common affine function \(1/4\), (7) is
\(\kappa_nf_r(z)\) on \([t-H_n,t+H_n]\).  Write
\(\mu_{D,P_n}(x)=E_{P_n}[D\mid X=x]\) for the observable conditional
treatment mean, equal to the appropriate side of (7).  For every measurable
prediction \(g\), conditional variance decomposition gives
\[
E_{P_n}[\{D-g(X)\}^2\mid X=x]
=\operatorname{Var}_{P_n}(D\mid X=x)
 +\{\mu_{D,P_n}(x)-g(x)\}^2.
\tag{12}
\]
The variance term in (12) is identical for the common-affine and
piecewise-affine models, so it cancels from their risk difference.
Let \(\mathcal A_{n,t}^0\) and \(\mathcal A_{n,t}^1\) denote,
respectively, the constrained common-affine and constrained
piecewise-affine classes used by the frozen construction at \(t\).  We
define, before using either object,
\[
 R_{n,t}(g)=E_{P_n}\!\left[
   \mathbf 1\{|X-t|\le H_n\}\{D-g(X)\}^2\right],
 \qquad
 Q_n(t)=\inf_{g\in\mathcal A_{n,t}^0}R_{n,t}(g)
       -\inf_{g\in\mathcal A_{n,t}^1}R_{n,t}(g).
\]

Here is the complete projection calculation.  With the basis
\(e(z)=(1,z)^\top\), the common, left-flank, and right-flank Gram
matrices and their inverses are
\[
 G_0=\int_{-1}^{1}e(z)e(z)^\top dz
   =\begin{pmatrix}2&0\\0&2/3\end{pmatrix},
 \qquad
 G_0^{-1}=\begin{pmatrix}1/2&0\\0&3/2\end{pmatrix},
 \tag{13}
\]
\[
 G_- =\int_{-1}^{0}e(z)e(z)^\top dz
   =\begin{pmatrix}1&-1/2\\-1/2&1/3\end{pmatrix},
 \qquad
 G_-^{-1}=\begin{pmatrix}4&6\\6&12\end{pmatrix},
 \tag{14}
\]
and
\[
 G_+ =\int_{0}^{1}e(z)e(z)^\top dz
   =\begin{pmatrix}1&1/2\\1/2&1/3\end{pmatrix},
 \qquad
 G_+^{-1}=\begin{pmatrix}4&-6\\-6&12\end{pmatrix}.
 \tag{15}
\]
The determinants are \(4/3\) for \(G_0\) and \(1/12\) for each of
\(G_-\) and \(G_+\).  They are therefore positive definite and full
rank, so every inversion in (13)--(15) is permitted.
If \(0\le r\le1\), the left moment is zero and both the common and
right moments equal
\[
 u(r)=\int_r^1 e(z)(1+z-r)\,dz
 =\begin{pmatrix}
 (r-3)(r-1)/2\\[2pt]
 (r-1)(r^2-2r-5)/6
 \end{pmatrix}.
 \tag{16}
\]
If \(-1\le r\le0\), the left, right, and common moments are,
respectively,
\[
 u_-(r)=\int_r^0 e(z)(1+z-r)\,dz
 =\begin{pmatrix}r(r-2)/2\\ r^2(r-3)/6\end{pmatrix},
 \tag{17}
\]
\[
 u_+(r)=\int_0^1 e(z)(1+z-r)\,dz
 =\begin{pmatrix}-(2r-3)/2\\ -(3r-5)/6\end{pmatrix},
 \qquad
 u_0(r)=u_-(r)+u_+(r)
 =\begin{pmatrix}
 (r-3)(r-1)/2\\[2pt]
 (r-1)(r^2-2r-5)/6
 \end{pmatrix}.
 \tag{18}
\]
For a regression function with moment vector \(u\) and Gram matrix
\(G\), the squared norm of its least-squares projection is
\(u^\top G^{-1}u\).  Let \(\overline Q_n(t)\) denote common risk minus
piecewise risk when the affine coefficient constraints are omitted.
After the change of variables \(x=t+H_nz\),
\[
 \overline Q_n(t)=H_n\kappa_n^2q(r),
 \tag{19}
\]
where (13)--(18) give, without an omitted remainder,
\[
 q(r)=
 \begin{cases}
 u_-(r)^\top G_-^{-1}u_-(r)
 +u_+(r)^\top G_+^{-1}u_+(r)
 -u_0(r)^\top G_0^{-1}u_0(r),&-1\le r\le0,\\[3pt]
 u(r)^\top G_+^{-1}u(r)-u(r)^\top G_0^{-1}u(r),&0\le r\le1,
 \end{cases}
 \tag{20}
\]
\[
 q(r)=
 \begin{cases}
 \displaystyle
 \frac{(1+r)^2(7r^4-32r^3+27r^2+22r+4)}{24},&-1\le r\le0,\\[7pt]
 \displaystyle
 \frac{(1-r)^2(7r^4-52r^3+99r^2-10r+4)}{24},&0\le r\le1.
 \end{cases}
 \tag{21}
\]
This also verifies directly that the population projection is the
object being evaluated, rather than assuming the desired profile
separation.  On the compact set of \(r\)'s used below, the coefficients
obtained from (13)--(18) are bounded numerical multiples of
\(\kappa_n\) in the \(z\) parametrization.  In the original
\(x\) parametrization, their slopes are therefore \(O(q_n)\).  Their
local-coordinate intercepts are \(1/4+O(\kappa_n)\); if an affine
function is instead written with a global \(x\)-intercept, that
intercept is \(1/4+O(q_n)\).  By (2), the coefficients in either
parametrization are strictly inside the compact constraints in
\(\mathrm{def:profile\mbox{-}cutoff\mbox{-}center}\) for all large
\(n\).  Hence \(Q_n(t)=\overline Q_n(t)\), and (19)--(21) hold with
\(Q_n\) in place of \(\overline Q_n\), at every point used in the
comparison.

The exact formula has a strict displacement gain:
\[
 q(0)=\frac16,
 \qquad
 q\left(\frac25\right)=\frac{5949}{31250},
 \qquad
 \Delta:=q\left(\frac25\right)-q(0)
       =\frac{1111}{46875}>0.
 \tag{22}
\]
Moreover, \(q(r)\le1/6\) for \(|r|\le1/8\).  On the negative half,
differentiating the first polynomial in (21) gives
\[
 q'(r)=\frac{(1+r)(7r^4-22r^3+2r^2+20r+5)}4>0,
 \qquad -\frac18\le r\le0,
 \tag{23}
\]
because the second factor is at least
\(5+20r\ge5/2\).  On the positive half, exact factorization gives
\[
 \frac16-q(r)
 =-\frac{r(r^2-4r+1)(7r^3-38r^2+51r-18)}{24}\ge0,
 \qquad 0\le r\le\frac18.
 \tag{24}
\]
Indeed, \(r^2-4r+1>0\) on this interval, whereas
\(7r^3-38r^2+51r-18\le7/512+51/8-18<0\).  This proves the asserted
local bound.

Let \(t_n^\circ=c-(2/5)H_n\).  This point is in \([1/4,3/4]\) because
\(H_n<1/4\).  The deterministic grid in the frozen construction has
mesh at most
\(\eta_n=\min\{n^{-1},m^{-1/3}H_n/64\}\), so it contains a point
\(\bar t_n\) such that
\[
 |\bar t_n-t_n^\circ|\le\eta_n,
 \qquad
 \frac{\eta_n}{H_n}\le\frac{m^{-1/3}}{64}\longrightarrow0.
 \tag{25}
\]
Continuity of the polynomial in (21), together with (22)--(25), yields
eventually
\[
 Q_n(\bar t_n)
 \ge H_n\kappa_n^2\left(\frac16+\frac{3\Delta}{4}\right),
 \qquad
 \sup_{|t-c|\le H_n/8}Q_n(t)
 \le\frac{H_n\kappa_n^2}{6}.
 \tag{26}
\]

It remains to show that random fitting cannot reverse the fixed gap in
(26).  We give the empirical comparison rather than assume a maximal
inequality.  Every localized squared-loss function used on the first
three folds is indexed by one interval center and at most four affine
coefficients in the fixed compact coefficient boxes of the
construction.  Since \(D\in\{0,1\}\), these losses have an envelope
\(B_M<\infty\) depending only on \(M\).  Put
\(\delta_m=m^{-3}\), take a deterministic \(\delta_m\)-mesh of the
cutoff-coordinate interval, and put a deterministic
\(\delta_m\)-net on every coefficient box.  There are at most
\(C_Mm^{C_M}\) representatives.  If \(t\) is replaced by its nearest
mesh point, every changed support or flank indicator lies in one of a
fixed collection of \(O(m^3)\) boundary bands, each having Lebesgue
length at most \(C\delta_m\).  Uniformity of \(X\) gives probability at
most \(C\delta_m\) to each band.  For a fold of size \(m\), write
\(N_B\) for its number of observations in a boundary band \(B\).  The
union bound and
\[
 P_n\{N_B\ge4\}\le {m\choose4}(C\delta_m)^4
\]
show that the probability some such band contains four observations
is \(O(m^3m^4\delta_m^4)=O(m^{-5})\).  Off this event, replacing \(t\)
by its mesh representative changes an empirical risk by
\(O_M(m^{-1})\).  If the affine slope is represented in the normalized
flank coordinate, its loss on an unchanged support varies by at most
\(O_M(\delta_m/H_n)\); in an unnormalized coordinate the smaller bound
\(O_M(\delta_m)\) applies.  Coefficient discretization contributes
\(O_M(\delta_m)\).  The corresponding population-risk errors satisfy
the same bounds, in addition to the \(O_M(\delta_m)\) boundary-mass
term, by the uniform design.  From (1),
\(m^{-3}/H_n=o(m^{-1})\).  Thus the deterministic finite net
approximates empirical and population risks uniformly by
\(O_M(m^{-1})\) on an event whose probability tends to one.

For completeness, the required finite-net concentration follows
directly from boundedness.  If \(W\) is centered and
\(|W|\le B_M\), convexity of the exponential on
\([-B_M,B_M]\) gives
\[
 E e^{aW}\le\cosh(aB_M)\le e^{a^2B_M^2/2}.
 \tag{27}
\]
The first inequality is the expectation of the chord bound for
\(e^{aw}\) on \([-B_M,B_M]\), using \(E[W]=0\).  The second follows
term by term from
\((2k)!\ge2^k k!\) in the power series for \(\cosh\).
Exponential Markov inequality applied to independent copies and then
optimized at \(a=t/B_M^2\) therefore gives
\[
 P\left(\left|m^{-1}\sum_{i=1}^mW_i\right|>t\right)
 \le2\exp\left\{-\frac{mt^2}{2B_M^2}\right\}.
 \tag{28}
\]
A union bound over the polynomial net, followed by the preceding
approximation, proves on each of the first three folds
\[
 \sup_{t,g}\left|
  \frac1m\sum_{i\in F_j}\mathbf 1\{|X_i-t|\le H_n\}
       \{D_i-g(X_i)\}^2
  -E_{P_n}\!\left[\mathbf 1\{|X-t|\le H_n\}
       \{D-g(X)\}^2\right]
 \right|
 =O_{P_n}\!\left(\sqrt{\frac{\log n}{n}}\right).
 \tag{29}
\]
The same construction with separate left- and right-flank indicators
covers the piecewise class.  Because the supremum in (29) ranges over
every admissible coefficient vector, it also covers a fitted function
trained on an independent fold.

Let \(e_n\) denote a common random upper bound in (29), defined over
the full deterministic loss classes.  Let \(\mathcal G_n\) be the
event that all empirical Gram checks in the frozen construction pass.
Only on \(\mathcal G_n\) are the fitted scores and their maximizer
formed.  On that event, uniform empirical-risk minimization implies,
for each model \(k\in\{0,1\}\), training fold \(j\), and candidate
\(t\),
\[
 0\le R_{n,t}(\widehat g_{j,t}^{k})
       -\inf_{g\in\mathcal A_{n,t}^k}R_{n,t}(g)\le2e_n.
 \tag{30}
\]
Applying (29) on the held-out fold and using (30) for the common and
piecewise fits gives, on \(\mathcal G_n\),
\[
 \sup_{t\in\mathcal T_n}|\widehat Q(t)-Q_n(t)|
 =O_{P_n}\!\left(\sqrt{\frac{\log n}{n}}\right).
 \tag{31}
\]
On \(\mathcal G_n^c\), the construction stops
before defining a score or a maximizer and returns
\(\mathsf F_{\mathrm{prof}}=1\); no undefined maximizer is invoked.
The empirical-process event underlying (29) is defined directly over
the full deterministic loss classes and hence is meaningful whether
or not \(\mathcal G_n\) occurs.

By (1),
\[
 \frac{\sqrt{\log n/n}}{H_n\kappa_n^2}
 =O\!\left(n^{-1/2+\varepsilon(2+1/s)}\sqrt{\log n}\right)
 \longrightarrow0.
 \tag{32}
\]
Let \(\mathcal E_n\) be the event, defined through the uniform
empirical-risk deviations over the deterministic classes, on which
those deviations are small enough that, whenever \(\mathcal G_n\)
occurs, the supremum in (31) is smaller than
\(H_n\kappa_n^2\Delta/4\).  Equations (29) and (32) give
\(P_n(\mathcal E_n)\to1\), whether or not the Gram checks pass.  On
\(\mathcal E_n\cap\mathcal G_n\), (26) shows that the score at
\(\bar t_n\) is strictly larger than the score at every grid point in
\([c-H_n/8,c+H_n/8]\).  The smallest maximizer \(\widetilde c\), which
is defined on \(\mathcal G_n\), must therefore satisfy
\[
 |\widetilde c-c|>H_n/8
 \quad\text{on }\mathcal E_n\cap\mathcal G_n.
 \tag{33}
\]
If \(\mathcal G_n^c\) occurs, the first probability in the target
display applies.  If \(\mathcal E_n\cap\mathcal G_n\) occurs, then the
later fitted-jump check either fails, in which case it again returns
\(\mathsf F_{\mathrm{prof}}=1\), or it passes, in which case
\(\widehat c_{\mathrm{prof}}=\widetilde c\) and (33) places the center
outside the stated radius.  The two events counted in the next display
are disjoint, and their union contains \(\mathcal E_n\).  Consequently,
\[
 P_n\{\mathsf F_{\mathrm{prof}}=1\}
 +P_n\!\left(\{|\widehat c_{\mathrm{prof}}-c|>H_n/8\}
       \cap\{\mathsf F_{\mathrm{prof}}=0\}\right)
 \ge P_n(\mathcal E_n)\longrightarrow1.
 \tag{34}
\]

Finally,
\[
 T_n=n\kappa_n^{2+1/s}
 =8^{-(2+1/s)}n^{\,1-\varepsilon(2+1/s)},
 \tag{35}
\]
so \(T_n/\log(e/\kappa_n)\to\infty\).  Also
\(nH_n\to\infty\), and therefore \(16/n\le H_n/8\) eventually.  For
arbitrary fixed positive
\(A_{\mathrm{prof}},a_{\mathrm{prof}},C_{\mathrm{prof}}\), the signal
condition with \(C_{\mathrm{prof}}\) eventually holds, whereas the
claimed right side at \(u_n=H_n/8\) is
\[
 A_{\mathrm{prof}}
 \exp\left\{-\frac{a_{\mathrm{prof}}m\kappa_n^2H_n}{8}\right\}
 =A_{\mathrm{prof}}
 \exp\left\{-\frac{a_{\mathrm{prof}}m}{8n(2^{12}M)^{1/s}}T_n\right\}
 \longrightarrow0,
 \tag{36}
\]
because \(m/n\to1/4\) and (35) diverges.  Equations (34) and (36)
contradict the frozen positive inequality for every choice of its three
constants and prove the stated negative result.

%%% FIELD smooth-bracket-replacement
For \(1<s\le2\), determine under the unchanged class
\(\mathcal P_n(s,M,v,\kappa_n)\) whether a genuinely different total
estimator or cutoff profile, together with same-class slope-kink lower
pairs, establishes the detectable-region frontier
\[
 L_n(\kappa_n)\asymp
 \min\left\{1,\frac{\rho_n}{\kappa_n}
                  +\frac1{n\kappa_n^3}\right\}
 \qquad\text{when}\qquad
 T_n\gtrsim\log(e/\kappa_n).
\]
The frozen common-versus-piecewise affine moving-window profile in the
\(1<s\le2\) branch of \(I_n^{\mathrm{SI}}\) cannot establish this upper
bracket through the proposed localization route:
the counterexample in
\(\mathrm{lem:profile\mbox{-}cutoff\mbox{-}localization}\) has
\(T_n/\log(e/\kappa_n)\to\infty\), yet the profile falls back or selects
a center farther than \(H_n/8\), where
\(H_n=(\kappa_n/(2^{12}M))^{1/s}\), with probability tending to one
because its population score is larger at a displaced split.  A positive
resolution must therefore construct and analyze a genuinely different
estimator or profile; alternatively, a negative resolution must prove
that the proposed frontier itself is unattainable.  Either resolution
must give explicit constants for uniform honesty and expected length,
derive a same-class lower construction for the term
\(1/(n\kappa_n^3)\), and decide whether oracle order is both sufficient
and necessary on the face \(n\kappa_n^2\rho_n\gtrsim1\).  No
cross-side derivative-matching assumption may be added, and the frozen
profile branch may not be reused as an attaining witness.

%%% FIELD smooth-bracket-partial
The known-cutoff theorem supplies the lower term
\(\rho_n/\kappa_n\), while
\(\mathrm{lem:profile\mbox{-}cutoff\mbox{-}localization}\) supplies an
explicit same-class obstruction to the frozen moving-window profile:
on its witness sequence the signal index dominates the logarithmic
search cost but the population score strictly favors a displaced split.
For a genuinely different cutoff estimator having a proved
high-probability radius \(R\), the fixed-center local-linear moment
calculation remains a conditional roadmap: the kernel identities
\(\int_0^1K_1(u)\,du=1\) and
\(\int_0^1uK_1(u)\,du=0\) yield a numerator radius of order
\(\rho_n+MR\), and division is valid because
\(d_P\ge\kappa_n>0\).  Thus a new estimator with
\(R\asymp(n\kappa_n^2)^{-1}\) would give the candidate order
\(\rho_n/\kappa_n+1/(n\kappa_n^3)\).  Neither such a replacement
estimator nor the matching slope-kink lower pair has yet been proved.

%%% FIELD critical-window-statement
Using \(\mathsf H_{\mathrm{crit}}\), determine a nonvariational
fixed-\(\alpha\) expression for \(L_n(\kappa_n)\) with a uniform
remainder throughout \(T_n\asymp\log(1/\kappa_n)\), and construct an
honest connected interval attaining that expression.  In particular,
for every \(\lambda\in(0,\infty)\), derive and match the limiting
expected length along sequences satisfying
\(T_n/\log(1/\kappa_n)\to\lambda\).  The result must retain
effect-label testing, agree with the exact \(s=1\) common-outcome
plateau experiment, have a \(\lambda\downarrow0\) converse that
recovers Theorem 2's nonshrinking lower bound, with no matched
below-boundary value claimed without a separate attaining upper bound,
and agree in order with Theorems 1 and 3 as
\(\lambda\uparrow\infty\).  It must strictly extend the known-cutoff
honesty result of Noack and Rothe (2024) by evaluating connected
expected length uniformly over the unknown cutoff \(c_P\).

%%% FIELD critical-window-partial
The explicit witness already evaluates one side of the critical
experiment.  Put \(C_h=(32/M)^{1/s}\) and
\(C_\psi=256/315\).  Then
\(h_n=C_h\kappa_n^{1/s}\),
\(\log J_n=s^{-1}\log(1/\kappa_n)+O(1)\), and
\[
 nC_\psi\kappa_n^2h_n=C_\psi C_hT_n.
\]
Thus, along \(T_n/\log(1/\kappa_n)\to\lambda\), the displayed mixture
chi-square divergence tends to zero whenever
\(C_\psi C_h\lambda<1/s\).  In that subregion the finite-sample
argument gives
\(\liminf_nL_n(\kappa_n)\ge2a_\star(1-2\alpha)\).  For \(s=1\), the
witness threshold is \(\lambda=315M/8192\).  At equality, the
first-order coordinate \(\lambda\) does not determine this mixture's
limit: the second-order value of
\(C_\psi C_hT_n-s^{-1}\log(1/\kappa_n)\) enters.  Above the witness
threshold its chi-square divergence grows and the construction
supplies no converse.  No honest connected interval with a matched
fixed-\(\alpha\) expected-length constant has been constructed in the
remaining critical region.
